% Options for packages loaded elsewhere
\PassOptionsToPackage{unicode}{hyperref}
\PassOptionsToPackage{hyphens}{url}
\documentclass[
  ignorenonframetext,
]{beamer}
\newif\ifbibliography
\usepackage{pgfpages}
\setbeamertemplate{caption}[numbered]
\setbeamertemplate{caption label separator}{: }
\setbeamercolor{caption name}{fg=normal text.fg}
\beamertemplatenavigationsymbolsempty
% remove section numbering
\setbeamertemplate{part page}{
  \centering
  \begin{beamercolorbox}[sep=16pt,center]{part title}
    \usebeamerfont{part title}\insertpart\par
  \end{beamercolorbox}
}
\setbeamertemplate{section page}{
  \centering
  \begin{beamercolorbox}[sep=12pt,center]{section title}
    \usebeamerfont{section title}\insertsection\par
  \end{beamercolorbox}
}
\setbeamertemplate{subsection page}{
  \centering
  \begin{beamercolorbox}[sep=8pt,center]{subsection title}
    \usebeamerfont{subsection title}\insertsubsection\par
  \end{beamercolorbox}
}
% Prevent slide breaks in the middle of a paragraph
\widowpenalties 1 10000
\raggedbottom
\AtBeginPart{
  \frame{\partpage}
}
\AtBeginSection{
  \ifbibliography
  \else
    \frame{\sectionpage}
  \fi
}
\AtBeginSubsection{
  \frame{\subsectionpage}
}
\usepackage{iftex}
\ifPDFTeX
  \usepackage[T1]{fontenc}
  \usepackage[utf8]{inputenc}
  \usepackage{textcomp} % provide euro and other symbols
\else % if luatex or xetex
  \usepackage{unicode-math} % this also loads fontspec
  \defaultfontfeatures{Scale=MatchLowercase}
  \defaultfontfeatures[\rmfamily]{Ligatures=TeX,Scale=1}
\fi
\usepackage{lmodern}
\ifPDFTeX\else
  % xetex/luatex font selection
\fi
% Use upquote if available, for straight quotes in verbatim environments
\IfFileExists{upquote.sty}{\usepackage{upquote}}{}
\IfFileExists{microtype.sty}{% use microtype if available
  \usepackage[]{microtype}
  \UseMicrotypeSet[protrusion]{basicmath} % disable protrusion for tt fonts
}{}
\makeatletter
\@ifundefined{KOMAClassName}{% if non-KOMA class
  \IfFileExists{parskip.sty}{%
    \usepackage{parskip}
  }{% else
    \setlength{\parindent}{0pt}
    \setlength{\parskip}{6pt plus 2pt minus 1pt}}
}{% if KOMA class
  \KOMAoptions{parskip=half}}
\makeatother
\setlength{\emergencystretch}{3em} % prevent overfull lines
\providecommand{\tightlist}{%
  \setlength{\itemsep}{0pt}\setlength{\parskip}{0pt}}
\usepackage{bookmark}
\IfFileExists{xurl.sty}{\usepackage{xurl}}{} % add URL line breaks if available
\urlstyle{same}
\hypersetup{
  hidelinks,
  pdfcreator={LaTeX via pandoc}}

\author{\texorpdfstring{}{}}
\date{}

\begin{document}

\begin{frame}
\newpage
\end{frame}

\begin{frame}[fragile]{The Frontier: Where We Need You}
\protect\phantomsection\label{sec:future}
We have built the foundation and validated the core mechanics. However,
the field is in its infancy. We have established valid trust metrics
(Trust Calculus) and filtering mechanisms (Firewalls), but we lack
standardized protocols for agent state persistence and identity
federation---the agentic equivalents of cookies or OAuth.

This is where you come in.

The Cognitive Integrity Framework is open source, and the problems below
are not just ``future research''---they are immediate engineering
blockers that need to be solved.

\begin{block}{1. The UX of Trust (Help Wanted)}
\protect\phantomsection\label{the-ux-of-trust-help-wanted}
\textbf{The Problem}: Currently, when an agent flag is raised
(``Identity Violation''), it looks like a JSON error log. \textbf{The
Need}: We need a ``Cognitive Dashboard'' for human operators. What does
it look like to visualize the trust graph of 50 agents in real-time? How
do we show ``Drift'' intuitively? \textbf{The Goal}: A React/Next.js
dashboard that connects to the CIF Python SDK.
\end{block}

\begin{block}{2. Standardized Agent Identity (Protocol Design)}
\protect\phantomsection\label{standardized-agent-identity-protocol-design}
\textbf{The Problem}: Every framework (LangChain, CrewAI, AutoGPT)
handles agent identity differently. \textbf{The Need}: A standard
``Agent Passport'' protocol. A cryptographically verifiable identity
token that an agent can carry across different frameworks. \textbf{The
Goal}: An RFC-style spec for \texttt{x-agent-identity} headers.
\end{block}

\begin{block}{3. ``Eusocial'' Security (Advanced Research)}
\protect\phantomsection\label{eusocial-security-advanced-research}
\textbf{The Problem}: Our current consensus is Byzantine, but it's
computationally expensive (\(O(n^2)\)). \textbf{The Need}: Insect
colonies don't vote; they use pheromones. We need \textbf{Stigmergic
Security Protocols} where agents leave ``trust trails'' in the
environment. \textbf{The Goal}: A lightweight, scalable consensus
algorithm modeled on ant colony immune responses.
\end{block}

\begin{block}{4. Benchmark Expansion}
\protect\phantomsection\label{benchmark-expansion}
\textbf{The Problem}: Our corpus has 950 attacks. Real-world capability
is growing daily. \textbf{The Need}: We need more attacks. Specifically,
we need \textbf{Multi-Modal Injection attacks} (audio/video) and
\textbf{Tool-Use Hijacking} examples. \textbf{The Goal}: Pull
Recommendations to the \texttt{cognitive\_integrity} repository adding
new scenarios to \texttt{data/attacks/}.
\end{block}
\end{frame}

\begin{frame}[fragile]
\begin{block}{Contributing}
\protect\phantomsection\label{contributing}
This is not a closed academic project. It is a living defense framework
for the agentic future.

\begin{itemize}
\tightlist
\item
  \textbf{Code}:
  \href{https://github.com/docxology/cognitive_integrity}{github.com/docxology/cognitive\_integrity}
\item
  \textbf{Discussion}: Join the \texttt{discussions} tab on GitHub.
\item
  \textbf{Contribute}: We prioritize PRs that add \textbf{Adapters} for
  new agent frameworks (e.g., Semantic Kernel, Microsoft AutoGen).
\end{itemize}

Let's build the immune system for the agentic web, together.
\end{block}
\end{frame}

\end{document}
